\documentclass[10pt]{article}
\usepackage[margin=0.78in]{geometry}
\usepackage{amsmath,amssymb,booktabs,array}
\usepackage[T1]{fontenc}
\usepackage[hidelinks]{hyperref}
\setlength{\parindent}{0pt}
\setlength{\parskip}{4pt}
\setlength{\emergencystretch}{2em}
\newcommand{\KL}{\operatorname{KL}}
\newcommand{\tr}{\operatorname{tr}}
\newcommand{\sg}{\operatorname{sg}}
\newcommand{\norm}[1]{\left\lVert#1\right\rVert}
\newcommand{\pos}[1]{\left[#1\right]_+}
\newcommand{\RR}{\mathbb R}
\title{JLZ v7: Joint Subject Targets Through a Causal Writer\\
\large Fresh actual keys, differentiable solves, and two allocation arms}
\author{ODE-edit method specification}
\date{October 2, 2026}
\begin{document}
\maketitle
\vspace{-1.5em}

\section{Scope and preserved native learning mechanism}
The method jointly learns subject-position interventions at every
eligible write layer and allocates them using the current model's actual
write geometry. Unlike a fixed-entry writer, an upper-layer basis is
recomputed after applying the candidate's lower-layer writes, and its
gradient is retained. This is a proposed design; this document does not
report production execution, GPU qualification, or measured improvement.

At batch entry freeze the model $W_t$, native anchor norms $a_{lr}>0$,
native KL teachers, and geometry statistic $A_l=\lambda_C C_{0l}+H_{tl}$.
Here $B$ is the actual logical batch size and $\mathcal L$ is the complete
ordered eligible layer set. Initialize absolute trainable columns
$D_l=[\delta_{l1},\ldots,\delta_{lB}]=0$. In the native virtual graph,
all requests use their usual native contexts, target tokens, lookup
positions, reductions, and readouts, with frozen $W_t$ and simultaneous
subject block-output interventions
\[
 h^{v}_{lrc}[s_{rc}]\leftarrow h^{v}_{lrc}[s_{rc}]+\delta_{lr}.
\]
The same request column is used in its rewrite contexts and native KL
input. Cross-layer propagation is in this common forward graph. No
official paraphrase, neighborhood example, future edit, or external
reference sentence enters optimization.

Define the mean native objective by
\begin{align}
 J_N(D)=\frac1B\sum_{r=1}^{B}\left[
 L^{v}_{N,r}(D)
 +\lambda_K\KL\bigl(p^v_{K,r}(D)\Vert p^0_{K,r}\bigr)
 +\lambda_n\sum_l\frac{\norm{\delta_{lr}}}{a_{lr}^2}\right].
 \label{eq:native}
\end{align}
NLL and KL retain their potentially different native readout layers.
The full-vocabulary native KL direction is current $\Vert$ entry.
The nonsquared native norm remains in the loss, with zero subgradient
chosen at zero. After each optimizer update, apply the individual native
clamp $\norm{\delta_{lr}}\leq c_{\rm native}a_{lr}$. All native
coefficients and anchors are recorded in the model profile.

Native task learning remains subject-mediated. A separate actual-weight
branch supplies differentiable geometry at every candidate and functional
feedback at four selected candidates. It uses no artificial subject
injection. Concentration in one layer is allowed; no layer is removed
before learning, and no entropy target or minimum layer quota is imposed.

\newpage
\section{Causal writer: fresh keys before each own-layer update}
For each candidate, traverse the actual branch in layer order. At layer
$l$, collect all current rewrite-context keys
$K_l^a(D_{<l})=[k^a_{lrc}]\in\RR^{d_{{\rm in},l}\times n}$ immediately
before that layer's modified projection. Lower-layer candidate weights
are already active. The current projection has not yet been modified for
this traversal. There is a layer barrier: collect the complete logical
batch's context keys before constructing that layer's update.

Let $\Omega=\operatorname{diag}(\alpha_{rc})$ with
$\alpha_{rc}\geq0$ and $\sum_c\alpha_{rc}=1$. Let
$Z\in\{0,1\}^{B\times n}$ indicate ownership of contexts, so that
$Z\Omega Z^\top=I_B$ under exact per-request normalization. Stored context
weights are preserved: casting FP32 weights to FP64 does not authorize
silent renormalization of their small summation error. In the initial
profile the canonical context has
weight $1/2$ and the noncanonical native contexts share the remaining
$1/2$. Abbreviating $K_l^a$ as $K_l$, compute
\begin{align}
 M_l&=A_l+K_l\Omega K_l^\top, &
 \bar K_l&=K_l\Omega Z^\top,\\
 P_l(D_{<l})&=\operatorname{solve}(M_l,\bar K_l), &
 U_l(D)&=D_lP_l^\top,\\
 W^a_l(D)&=W_{t,l}^{32}
       +\operatorname{cast}_{32}\!\left(D_l^{64}(P_l^{64})^\top\right).
 \label{eq:causal}
\end{align}
Apply this weight to all token positions, continue to the next layer,
and repeat. Every candidate starts from $W_t$, not from the preceding
candidate's edited model. All current requests determine each basis,
including when functional feedback later uses a subset.

For fixed actual keys, Eq.~\eqref{eq:causal} minimizes
\[
 \frac12\norm{(UK_l-D_lZ)\Omega^{1/2}}_F^2
 +\frac12\tr(UA_lU^\top).
\]
Here $\delta$ is an incremental own-layer target: $UK_l\approx D_lZ$.
It is not an absolute target subsequently corrected by subtracting a
newly measured activation. In particular,
\[
 h^a_l-h^v_l=(h^{a,\mathrm{pre}}_l-h^{v,\mathrm{pre}}_l)
              +(U_lK_l-D_lZ)
\]
at the corresponding context columns and supported block-output sites.
Same-layer fitting alone need not remove inherited lower-layer error.
Replacing $D_lZ$ by an absolute virtual-state residual would define a
different writer and is not part of this version.

The first writable layer's keys and basis are invariant within a batch,
because their computation precedes every modified projection. Cache
them once. Every upper-layer basis is fresh for each candidate and
retains gradients through keys and solves. The construction is causal,
not a simultaneous fixed-point iteration: $P_l$ depends on $D_{<l}$ but
not on an own-layer key produced after its own write.
An adapter must verify independent write parameters and this feed-forward
dependency. Aliased weights, recurrent invocations, or shared write sites
are not supported merely by naming corresponding layers.

\newpage
\section{Differentiable solve and exact reuse boundaries}
Factor each frozen positive-definite $A_l=L_lL_l^\top$ once per batch.
A context-dual implementation can compute
\begin{align}
 F_l&=L_l^{-1}K_l\Omega^{1/2}, & S_l&=I+F_l^\top F_l,\\
 V_l&=\operatorname{solve}(S_l,\Omega^{1/2}Z^\top), &
 T_l&=F_lV_l, & P_l&=L_l^{-\top}T_l.
 \label{eq:dual}
\end{align}
This is algebraically equivalent to the primal solve. It uses no
division by context weights and permits zero weights. Geometry and
solves use FP64. Shape, factorization, and residual checks may detect
numerical failure; silent diagonal jitter changes the writer and is not
an approved recovery. Neither contexts nor off-diagonal terms are dropped.

The differential for the solve subgraph, with $A,\Omega,Z$ fixed, is
\[
 dP=M^{-1}\!\left[dK\Omega(Z^\top-K^\top P)
                      -K\Omega\,dK^\top P\right].
\]
Given upstream adjoint $Q=\nabla_P L$, let
$Y=\operatorname{solve}(M^\top,Q)$. Its contribution is
\[
 \nabla_K L\big|_{\rm solve}
 =Y(Z-P^\top K)\Omega-P(Y^\top K)\Omega.
 \label{eq:vjp}
\]
Direct key dependencies from geometry or activations are added separately;
Eq.~\eqref{eq:vjp} is not the entire loss gradient. For upper layers,
\[
 dU_l=dD_lP_l^\top+D_l\,dP_l^\top
\]
couples lower-layer interventions to later-layer write geometry. Detaching
$K$, reusing a preceding candidate's $P$, or omitting the solve adjoint
removes this coupling and defines another method.

For a linear operation with input $X$, output adjoint $G$, and update
$DP^\top$, the real-arithmetic identities
\[
 \nabla_D L=G^\top(XP),\qquad
 \nabla_P L=X^\top(GD),\qquad
 \nabla_X L=GW^a
\]
can avoid materializing a full weight-gradient tensor. Both parameter
adjoints and the full input adjoint are required. Forward
materialization, casts, and rounding remain as declared. A reordered
custom backward must be validated against the ordinary autograd
reference; these identities do not imply bitwise floating-point equality.

Checkpointing may recompute exact candidate operations. Its closure must
be a pure function of immutable candidate weights/state or install those
weights inside recomputation. It cannot depend on a temporary hook
context that expired after forward. Native and actual graphs have
different cache boundaries. Upper-layer key activations and factors
cannot be cached across optimizer candidates.

\newpage
\section{Dynamic geometry costs and the two arms}
At every candidate, construct the differentiable small matrices
\begin{align}
 G_l(D_{<l})&=P_l^\top A_lP_l=T_l^\top T_l,\\
 E_l(D_{<l})&=(P_l^\top K_l-Z)\Omega(P_l^\top K_l-Z)^\top,\\
 G_l+E_l&=Z\Omega Z^\top-\bar K_l^\top P_l.
 \label{eq:grams}
\end{align}
The last identity uses the actual stored weights; its first term is
$I_B$ in exact normalized arithmetic. It is a numerical check, not a
reason to ignore cancellation in small residual energies. With the frozen
$\sigma_l=(B^{-1}\sum_r a_{lr}^2)^{1/2}$, define
\[
 g_l(D)=\frac{\sqrt{\tr(D_lG_l(D_{<l})D_l^\top)}}{\sqrt B\,\sigma_l},
 \qquad
 e_l(D)=\frac{\sqrt{\tr(D_lE_l(D_{<l})D_l^\top)}}{\sqrt B\,\sigma_l}.
\]
These correspond to actual-key weighted write action and same-layer
incremental realization error. Explicit factors such as
$D_l(P_l^\top K_l-Z)\Omega^{1/2}$ permit stable norm evaluation.
These energies refer to the FP64 theoretical update; telemetry separately
measures the effects of the materialized FP32 candidate.
All derivatives through $G,E,P,K$ are included; an eigendecomposition
followed by detached factors is not an equivalent implementation.
Use zero subgradient at an exact zero norm, with no root-norm smoothing.
Adding epsilon smoothing requires a separate method revision.

\begin{center}
\begin{tabular}{lll}
\toprule
Arm & $R_W(D)$ & $R_E(D)$\\
\midrule
A: concentration permitted & $\sum_lg_l$ & $\sum_le_l$\\
B: spread preference & $\sqrt{\sum_lg_l^2}$ & $\sqrt{\sum_le_l^2}$\\
\bottomrule
\end{tabular}
\end{center}
Both arms use the primary mean objective
\[
 J_{\rm base}^{(a)}=J_N+\lambda_WR_W^{(a)}+\lambda_ER_E^{(a)},
 \qquad \lambda_W=\lambda_E=0.1
\]
in the initial, unvalidated profile. For fixed nonnegative cost entries
of fixed total, the Euclidean aggregation discounts spreading relative
to a sum. This describes only the outer aggregation. Because the
geometry itself now depends on lower-layer parameters, these are not
globally convex norms of the full collection $D$ and the same allocation
change can also change all subsequent costs. Arm A does not make the
complete objective allocation-neutral. Common hidden-feedback penalties
can also favor particular distributions.

Both arms consider all layers and share the native loss, physical
feedback, memory, Adam settings, clamp, and evaluation. They differ only
in the two outer cost aggregations. Concentration is a possible learned
outcome, not an initial layer restriction or an acceptance condition.
Geometry costs are preservation proxies; they do not establish functional
locality or rule out an under-edited optimum.

\newpage
\section{Scheduled physical feedback from native sentences}
At pulse candidates 5, 10, 15, and 20, use the freshly built causal
weights as the physical student. The teacher is the same candidate's
native jointly injected forward, detached only on the teacher side.
For a current cohort $I$, native rewrite weights $w_{rc}$, and target
lengths $T_{rc}$, define
\begin{align}
 C_h(I)&=\frac1{|I|}\sum_{r\in I}\frac12\sum_l\sum_cw_{rc}
  \frac{\norm{h^a_{lrc}-\sg(h^v_{lrc})}^2}{a_{lr}^2},\\
 C_d(I)&=\frac1{|I|}\sum_{r\in I}\sum_cw_{rc}\frac1{T_{rc}}
  \sum_{j=1}^{T_{rc}}\KL\bigl(\sg(p^v_{rcj})\Vert p^a_{rcj}\bigr),\\
 C_{\rm cur}(I)&=0.1C_h(I)+0.1C_d(I)
  +\frac{\lambda_K}{|I|}\sum_{r\in I}\KL(p^a_{K,r}\Vert p^0_{K,r}).
\end{align}
Hidden states are total subject block outputs after intervention.
Distillation uses full-vocabulary probabilities at native NLL target
positions and readout; physical KL uses the native KL input and readout.
Teacher $\Vert$ student distillation and current $\Vert$ entry native
KL deliberately have different directions. No additional physical target
NLL replaces the virtual native task objective.

For eligible past native inputs $J$, use the same physical weights and
the same loss in both arms:
\[
 C_{\rm past}(J)=\frac1{|J|}\sum_{r\in J}\left[
  \lambda_K\KL(p^a_{K,r}\Vert p^{\rm admit}_{K,r})
  +\sum_cw_{rc}\frac12\pos{L^a_{N,rc}-b^{\rm commit}_{rc}}^2\right].
\]
The admission-entry KL teacher stays fixed while its identity is
resident. A retained repeated or readmitted fact uses its latest target
and actual committed NLL baseline. Memory follows the native-only v5
unique-fact Algorithm R reservoir of capacity 128, including shared-token
teacher ownership and eviction. Current fact identities are excluded
from past preservation; admission has no quality gate.

Once per batch, seed a disjoint partition $I_1,\ldots,I_4$ covering all
$B$ current requests. Select $M=\min(16,B,\text{eligible memory size})$
past requests once and partition them into disjoint $J_1,\ldots,J_4$.
At candidate $q=5j$ add
\[
 \frac{|I_j|}{B}C_{\rm cur}(I_j)+\frac{|J_j|}{M}C_{\rm past}(J_j)
\]
to $J_{\rm base}^{(a)}$, defining $J_q$. Empty terms are zero, including
all past terms when $M=0$. Cohort weights sum to one across pulses, with
no factor-four multiplier. Every selected request contributes its full
native bundle. Other candidates have no functional-feedback term but
still build actual keys and differentiate dynamic geometry on all current
rewrite inputs. Moving detached teachers make this a specified finite
schedule rather than descent on a single fixed global objective.

\newpage
\section{Adam, exact final commit, and sequential state}
Use absolute-coordinate Adam with
$(\beta_1,\beta_2)=(0.9,0.999)$, $\epsilon=10^{-8}$, weight decay zero,
and fresh moments at batch entry. Candidates $q=1,\ldots,25$ evaluate
$D^{q-1}$. Candidates 1--24 accumulate the complete logical gradient
\[
 \nabla_D\{B J_q(D^{q-1})\}
\]
and perform one update with constant native-profile learning rate
$\gamma_q=\gamma_{\rm native}$; the initial profile uses $0.1$ from
update 1. There is no warmup. Multiplication by $B$ applies to the entire
mean loss, including pulse terms. Microbatching changes neither this
normalization nor the number of optimizer steps. Adam epsilon makes
silent gradient rescaling a material implementation change.

Apply the individual native clamp after every update. Clamp activity
does not change learning rate or introduce a quality gate. There is no
line search or early stop based on task quality. Nonfinite values,
invalid solves, and state/shape mismatches remain numerical failures,
not grounds to silently remove layers or change coefficients.

Candidate 25 has no backward or update. Rebuild its causal writer in
the same layer order, evaluate the complete current actual native bundles,
and commit precisely those materialized weight tensors. There is no
alternative final writer, gain multiplier, residual redefinition, or
post-hoc re-solve with a different objective. Record equality between
the physical tensors evaluated and the tensors committed.
Re-capture the terminal physical pre-own-layer keys and verify their
agreement with the builder keys: later causal writes cannot change
earlier keys. Final inference has no subject-intervention hook.

The final traversal provides actual full-context keys. Append once per
encountered request to history:
\[
 H_{t+1,l}=H_{tl}+\sum_{r,c}\alpha_{rc}
 k^{a,\rm final}_{lrc}(k^{a,\rm final}_{lrc})^\top.
\]
Do not append history at intermediate candidates or each pulse. Update
memory with actual final per-context NLL and the prescribed admission
teacher. Frozen historical key statistics can still become stale across
later batches; refreshing current candidate keys does not solve that
separate historical approximation.

Record virtual and physical native NLL/KL, native norm, dynamic geometry
costs and gradients, per-layer payload norms, actual write norms,
realization error, clamp activity, pulse memberships, solve residuals,
and evaluated-to-committed equality. Evaluation R/P/N, including cohort
at-write and cumulative denominators, remains observational and never
feeds optimization. These diagnostics are not layer-selection gates.
Terminal gradients are unmeasured and recorded as null, not as zero.

The method accepts arbitrary positive logical $B$ and compatible
model/benchmark profiles. An adapter must identify the writable projection,
corresponding subject block-output site, and exact native readouts. The
first experiment is two arms with batch size 100 for five sequential
batches, 500 edits per arm; these limits do not define the method.
The initial profile is Llama3-8B-Instruct/CounterFact, layers L4--L8,
$\lambda_n=0.5$, $\lambda_K=0.0625$, $c_{\rm native}=0.75$,
and $\lambda_C=15000$, with FP32 model weights and FP64 geometry.
Each arm starts from its own cold model/history state.

\newpage
\section{Cost, qualification, and unresolved scientific questions}
The fixed-writer v6 exposure count and any proposed percentage speedup
cannot be carried over to v7. In addition to 25 native candidate forward
evaluations and 24 native backwards, the reference schedule requires
25 complete-current-rewrite key-builder traversals and 24 corresponding
differentiable traversals. Upper-layer factorization/solve work occurs
at every candidate, with adjoints on the first 24. First-layer keys/basis
and the frozen $A_l$ factorizations are computed once per batch.

\begin{center}
\begin{tabular}{p{0.48\linewidth}p{0.41\linewidth}}
\toprule
Component & Reference frequency\\
\midrule
Native joint subject candidate & 25 forward, 24 backward\\
Full-current rewrite key builder & 25 traversals, 24 with gradient\\
Upper-layer dynamic basis & $25(|\mathcal L|-1)$ fresh solves and
$24(|\mathcal L|-1)$ solve adjoints\\
Physical feedback cohorts & Four pulses; current requests covered once,
past $M$ requests covered once\\
Terminal physical current audit & All current native bundles once,
sharing the final candidate's causal weights\\
\bottomrule
\end{tabular}
\end{center}
These counts are not equal-cost full-model forwards or independent
additive FLOP estimates. Key-builder prefixes can supply current pulse
activations; physical KL/past rows and readout suffixes have different
costs. Instrument actual forward/backward operations, processed tokens,
factorizations, wall time, and peak memory. An optional zero-candidate
identity fast path requires equivalence checks and separate accounting;
the reference count remains 25 builder traversals. Small physical cohorts
do not remove the all-current builder dependency.
For geometry alone the builder may stop at the last writable site's
input key; constructing its update does not require the model's final
readout suffix. Full physical observations still apply that final write.

Allowed optimizations preserve the graph: exact checkpointing; shared
candidate materialization; immutable prefix caches; selected-position
full-vocabulary heads; and validated implicit/custom VJPs. Detaching
current keys or solve gradients, reusing upper factors across candidates,
or restricting the builder to pulse requests changes the method. No
production/GPU run is claimed by this specification.

A causally qualified adapter may crop only the key-builder rows after
each subject while preserving original tokens, masks, positions, and
all contexts. The reference builder uses full rows. Native NLL inputs
remain complete, and current physical feedback plus the terminal audit
still require complete prompts, unless equivalence of a cached
continuation is separately qualified. Prefix cropping does not remove
the requirement to include every current request in the writer solve.

Qualification must compare primal and dual solves, check the Gram
identity, finite-difference or autograd-check solve and total causal
gradients, validate the direct-linear VJP under declared dtypes, test
cross-request and lower-to-upper paths, and verify identical evaluated
and committed tensors. Include $B<4$, empty past memory, zero context
weights, the zero initialization, and checkpoint recomputation after
the original forward context has exited.

The stronger connection is precise but limited: the native subject target
now receives gradient through how its lower-layer write changes later
keys and writer feasibility. This does not make all targets realizable,
guarantee paraphrase generalization, or cover all neighboring knowledge.
The provisional coefficients, finite Adam budget, native-only memory,
and sparse functional pulses remain empirical design choices. Arm A/B
isolates the outer allocation-cost aggregation; other changes require
separate comparisons. No speed or quality claim is made before
qualification and measured experiments.
\end{document}
